% Bagian: computability
% Bab: computability-theory
% Subbagian: no-universal-function

\documentclass[../../../include/open-logic-section]{subfiles}

\begin{document}

\olfileid[id]{cmp}{thy}{nou}
\olsection{Tidak Ada Fungsi Dapat Dihitung Universal}

Meskipun terdapat suatu fungsi dapat dihitung parsial yang universal bagi
fungsi-fungsi dapat dihitung parsial, tidak ada fungsi dapat dihitung total yang
universal bagi fungsi-fungsi dapat dihitung total.

\begin{thm}
\ollabel{thm:no-univ}
Tidak ada fungsi dapat dihitung universal. Dengan kata lain, setiap fungsi
$\fn{Un}'(k, x)$ yang memenuhi sifat bahwa jika $f(x)$ merupakan fungsi dapat
dihitung total, maka terdapat suatu bilangan asli~$k$ sedemikian sehingga
$f(x) = \fn{Un}'(k,x)$ untuk setiap~$x$, bukanlah fungsi yang dapat dihitung.
\end{thm}

\begin{proof}
Buktinya merupakan diagonalisasi sederhana: jika $\fn{Un}'(k,x)$ bersifat
total dan dapat dihitung, maka
\[
d(x) = \fn{Un}'(x, x) + 1
\]
juga bersifat total dan dapat dihitung. Namun, berdasarkan definisi, $d(k)$
tidak sama dengan $\fn{Un}'(k,k)$. Jadi, untuk setiap $k$, nilai $d(x)$ dan
$\fn{Un}'(k, x)$ berbeda untuk sekurang-kurangnya satu~$x$, yaitu $x = k$.
\end{proof}

\begin{explain}
\olref[uni]{thm:univ-comp} di atas menunjukkan bahwa kita dapat menghindari
argumen diagonalisasi ini, tetapi hanya dengan mengizinkan fungsi universal
bersifat parsial. Artinya, $\fn{Un}$ bersifat universal bagi fungsi-fungsi dapat
dihitung total, hanya saja fungsi itu sendiri tidak total. Argumen
diagonalisasi tidak berlaku dalam kasus parsial.
\end{explain}

\begin{prob}
  Untuk memahami mengapa argumen diagonalisasi dalam bukti
  \olref{thm:no-univ} tidak berlaku dalam kasus parsial, tinjau fungsi
  $f(x) \simeq \fn{Un}(x,x)+1$. Apakah fungsi itu dapat dihitung secara
  parsial? Jika demikian, fungsi itu mempunyai indeks~$e$, yakni
  $f(x) \simeq \fn{Un}(e,x)$. Apa yang dapat Anda katakan tentang~$f(e)$?
\end{prob}


\end{document}

